% =====================================================================
% પરિશિષ્ટો : સ્થિરાંકો, સૂત્ર-સંગ્રહ, ઉત્તર-કી, પરિભાષા કોશ, સંદર્ભગ્રંથો
% =====================================================================
\appendix

% ---------------------------------------------------------------------
\chapter[પરિશિષ્ટ A : સ્થિરાંકો]{મૂળભૂત ભૌતિક સ્થિરાંકો\\{\large (Fundamental Physical Constants)}}
\label{app:constants}

\begin{center}
\begin{tabular}{@{}llll@{}}
\toprule
\textbf{સ્થિરાંક} & \textbf{ચિહ્ન} & \textbf{મૂલ્ય} & \textbf{એકમ} \\
\midrule
પ્રકાશની ઝડપ (Speed of Light)          & $c$      & $2.998\times10^{8}$        & m\,s$^{-1}$ \\
પ્લાંક સ્થિરાંક (Planck Constant)       & $h$      & $6.626\times10^{-34}$      & J\,s \\
ઘટાયેલ પ્લાંક સ્થિરાંક                  & $\hbar$  & $1.055\times10^{-34}$      & J\,s \\
વિદ્યુત વીજભાર (Elementary Charge)       & $e$      & $1.602\times10^{-19}$      & C \\
ઇલેક્ટ્રોન-સ્કંધ (Electron Mass)         & $m_e$    & $9.109\times10^{-31}$      & kg \\
પ્રોટોન-સ્કંધ (Proton Mass)              & $m_p$    & $1.673\times10^{-27}$      & kg \\
ન્યૂટ્રોન-સ્કંધ (Neutron Mass)           & $m_n$    & $1.675\times10^{-27}$      & kg \\
બોલ્ટ્ઝમાન સ્થિરાંક                      & $k_B$    & $1.381\times10^{-23}$      & J\,K$^{-1}$ \\
શૂન્યાવકાશ પારગમ્યતા                     & $\varepsilon_0$ & $8.854\times10^{-12}$ & F\,m$^{-1}$ \\
શૂન્યાવકાશ ભેદ્યતા                        & $\mu_0$  & $4\pi\times10^{-7}$        & H\,m$^{-1}$ \\
રીડબર્ગ ઊર્જા (Rydberg Energy)          & $R_\infty hc$ & $13.606$             & eV \\
બોર ત્રિજ્યા (Bohr Radius)               & $a_0$    & $5.292\times10^{-11}$      & m \\
ઇલેક્ટ્રોન કૉમ્પ્ટન તરંગલંબાઈ            & $\Lambda_e$ & $2.426\times10^{-12}$   & m \\
સ્ટેફન-બોલ્ટ્ઝમાન સ્થિરાંક                & $\sigma$ & $5.670\times10^{-8}$       & W\,m$^{-2}$K$^{-4}$ \\
વિન વિસ્થાપન સ્થિરાંક                    & $b$      & $2.898\times10^{-3}$       & m\,K \\
અણુભાર એકમ                              & $u$      & $1.6605\times10^{-27}$     & kg \\
\bottomrule
\end{tabular}
\end{center}

\paragraph{ઉપયોગી સંખ્યાત્મક સંબંધો:}
\begin{align*}
  E_\gamma(\mathrm{eV}) &= \frac{1240}{\lambda(\mathrm{nm})}
  & \lambda(\text{\AA}) &= \frac{12.27}{\sqrt{V(\mathrm{volt})}}\quad(\text{ઇલેક્ટ્રોન})\\[1mm]
  \kB T(T=300\,\mathrm{K}) &= 25.9\ \mathrm{meV}
  & \Delta\lambda_C(180^\circ) &= 4.86\ \mathrm{pm}
\end{align*}

% ---------------------------------------------------------------------
\chapter[પરિશિષ્ટ B : સૂત્ર-સંગ્રહ]{સૂત્ર-સંગ્રહ : પરીક્ષા-તૈયારી પત્રિકા\\{\large (Formula Sheet for GATE/NET Revision)}}

\section*{એકમ 1 -- જૂનો ક્વોન્ટમ વાદ}
\begin{itemize}
  \item પ્લાંક: $u(\nu,T)=\dfrac{8\pi h\nu^{3}}{c^{3}}\dfrac{1}{e^{h\nu/k_BT}-1}$;\quad
        $\bar\varepsilon=\dfrac{h\nu}{e^{h\nu/k_BT}-1}$
  \item સ્ટેફન-બોલ્ટ્ઝમાન $R=\sigma T^{4}$;\quad વિન $\lambda_{\max}T=b$
  \item આઇન્સ્ટાઇન $eV_s=h\nu-\phi=h(\nu-\nu_0)$
  \item કૉમ્પ્ટન $\Delta\lambda=\dfrac{h}{m_ec}(1-\cos\theta)=2\Lambda\sin^{2}\dfrac{\theta}{2}$
  \item ડી બ્રોગ્લી $\lambda=h/p$;\quad ઔષ્ધિજ $\lambda=h/\sqrt{3mk_BT}$
  \item બોર-સોમરફેલ્ડ $\oint p_i\,\dd q_i=n_i h$
\end{itemize}

\section*{એકમ 2 -- દ્વૈતતા અને અનિશ્ચિતતા}
\begin{itemize}
  \item $v_p=\omega/k$;\quad $v_g=\dd\omega/\dd k=v_p-\lambda\,\dd v_p/\dd\lambda$;
        મુક્ત કણ: $v_g=v$, $v_pv_g=c^2$ (સાપેક્ષિક)
  \item પેકેટ-ફેલાવો $\sigma(t)=\sigma\sqrt{1+(\hbar t/m\sigma^{2})^{2}}$
  \item $\Delta x\,\Delta p_x\ge\dfrac{\hbar}{2}$;\quad
        $\Delta E\,\Delta t\ge\dfrac{\hbar}{2}$;\quad
        $\Delta L_z\,\Delta\varphi\ge\dfrac{\hbar}{2}$
  \item કુદરતી પહોળાઈ $\Gamma=\dfrac{\hbar}{\tau}$;\quad શૂન્ય-બિંદુ ઊર્જા $E_0=\dfrac{\hbar\omega}{2}$
\end{itemize}

\section*{એકમ 3--4 -- શ્રોડિન્જર યાંત્રિકી અને ચોક્કસ ઉકેલો}
\begin{itemize}
  \item TDSE: $i\hbar\dfrac{\partial\Psi}{\partial t}=\hat H\Psi$;\quad
        TISE: $-\dfrac{\hbar^{2}}{2m}\nabla^{2}\psi+V\psi=E\psi$
  \item અપેક્ષિત મૂલ્ય $\langle A\rangle=\displaystyle\int\Psi^{*}\hat A\Psi\,\dd^3r$
  \item અનંત કૂવો: $E_n=\dfrac{n^{2}\pi^{2}\hbar^{2}}{2mL^{2}}$;\quad
        $\psi_n=\sqrt{\dfrac{2}{L}}\sin\dfrac{n\pi x}{L}$
  \item આવર્તક દોલક: $E_n=(n+\tfrac12)\hbar\omega$, Hermite બહુપદી ઉકેલ
  \item અવરોધ-સુરંગન (Tunneling): $T\simeq\dfrac{16E(V_0-E)}{V_0^{2}}\,e^{-2\kappa a}$,
        $\kappa=\sqrt{2m(V_0-E)}/\hbar$; ચોક્કસ સ્વરૂપ એકમ 4 માં
  \item Ehrenfest: $\langle x\rangle$ ચિરસંમત ગતિ-સમીકરણને અનુસરે છે
\end{itemize}

\section*{એકમ 5 -- ક્ષોભ સિદ્ધાંત, આંકડાશાસ્ત્ર અને આધુનિક તકનીકો}
\begin{itemize}
  \item ક્ષોભ સિદ્ધાંત: $E_n^{(1)}=\langle n^{(0)}|\hat H'|n^{(0)}\rangle$
  \item ફર્મી ઊર્જા $E_F=\dfrac{\hbar^{2}}{2m}(3\pi^{2}n)^{2/3}$
  \item FD/BE: $f(E)=\bigl[e^{(E-\mu)/k_BT}\pm1\bigr]^{-1}$
  \item આઇન્સ્ટાઇન સંબંધ $A_{21}/B_{21}=8\pi h\nu^{3}/c^{3}$, $B_{12}=B_{21}$
  \item ક્યુબિટ $|\psi\rangle=\alpha|0\rangle+\beta|1\rangle$, $|\alpha|^{2}+|\beta|^{2}=1$
\end{itemize}

% ---------------------------------------------------------------------
\chapter[પરિશિષ્ટ C : MCQ ઉત્તર-કી]{MCQ ઉત્તર-કી સંકલન\\{\large (Consolidated Answer Key)}}
\label{app:answers}

દરેક એકમની MCQ-યાદીની ઉત્તર-કી ત્યાં જ `javabchavi` બોક્સમાં આપેલ છે. સગવડ માટે એકમ 1 અને 2 ની
સંક્ષિપ્ત કી અહીં પુનઃ આપાય છે.

\section*{એકમ 1}
\begin{center}
\begin{tabular}{@{}l*{10}{c}@{}}
\toprule
પ્રશ્ન & 1 & 2 & 3 & 4 & 5 & 6 & 7 & 8 & 9 & 10\\
\midrule
ઉત્તર & b & d & a & b & a & b & a & a & b & b\\
\bottomrule
\end{tabular}
\end{center}

\section*{એકમ 2 (50 MCQ)}
પૂર્ણ કી અને ગણતરી-આધારિત પ્રશ્નોની પગલાંવાર સમજૂતી એકમ 2 ના અંતે આપેલ છે
(1-B, 2-C, 3-B, 4-B, 5-C, \dots, 50-C).

\paragraph{સ્વ-પરીક્ષણ સૂચન:} કોઈ પ્રશ્નમાં ભૂલ હોય તો સંબંધિત એકમના વિભાગ પર પાછા જાઓ --
દરેક MCQ પુસ્તકના કોઈક સૂત્ર અથવા ઉત્પત્તિ સાથે સીધો સંબંધ ધરે છે.

% ---------------------------------------------------------------------
\chapter[પરિશિષ્ટ D : પરિભાષા કોશ]{દ્વિભાષી પરિભાષા કોશ\\{\large (Gujarati--English Glossary of Technical Terms)}}
\label{app:glossary}

\begin{longtable}{@{}p{0.42\linewidth}p{0.52\linewidth}@{}}
\toprule
\textbf{ગુજરાતી પરિભાષા} & \textbf{English Term} \\
\midrule
\endfirsthead
\toprule
\textbf{ગુજરાતી પરિભાષા} & \textbf{English Term} \\
\midrule
\endhead
કૃષ્ણિકા વિકિરણ & Black body radiation \\
વર્ણપટ ઊર્જા ઘનતા & Spectral energy density \\
અલ્ટ્રાવાયોલેટ આપત્તિ & Ultraviolet catastrophe \\
મોડ ઘનતા & Mode density \\
પ્રકાશ-વિદ્યુત અસર & Photoelectric effect \\
કાર્ય વિધેય & Work function \\
દેહલી આવૃત્તિ & Threshold frequency \\
અટકાવતી વોલ્ટેજ & Stopping potential \\
ફોટોન / પ્રકાશ ક્વોન્ટમ & Photon / Light quantum \\
કૉમ્પ્ટન વિક્ષેપણ & Compton scattering \\
પદાર્થ તરંગ & Matter wave \\
તરંગ-કણ દ્વૈતતા & Wave--particle duality \\
કળા ઝડપ / સમૂહ ઝડપ & Phase velocity / Group velocity \\
તરંગ પેકેટ & Wave packet \\
ડિસ્પર્શન & Dispersion \\
અનિશ્ચિતતા સિદ્ધાંત & Uncertainty principle \\
વિચાર પ્રયોગ & Gedanken experiment \\
શૂન્ય-બિંદુ ઊર્જા & Zero-point energy \\
કુદરતી રેખા-પહોળાઈ & Natural line broadening \\
તરંગ વિધેય & Wave function \\
જન્મ-સંભાવના અર્થઘટન & Born probability interpretation \\
સંભાવ્યતા ધારા ઘનતા & Probability current density \\
ઑપરેટર / સંક્રિયક & Operator \\
હર્મિટિયન (સ્વયં-સહયોગી) ઑપરેટર & Hermitian operator \\
આગેક્રમ (કમ્યુટેટર) & Commutator \\
અપેક્ષિત મૂલ્ય & Expectation value \\
પ્રમાણિત & Normalized \\
સ્થિતિ-ઊર્જા કૂવો & Potential well \\
સ્થિતિ-ઊર્જા અવરોધ & Potential barrier \\
સુરંગન અસર & Tunneling effect \\
આવર્તક દોલક & Harmonic oscillator \\
કોણીય સંવેગ & Angular momentum \\
ગોળીય હાર્મોનિક વિધેયો & Spherical harmonics \\
કક્ષીય / ચુંબકીય સંખ્યા & Orbital / Magnetic quantum number \\
ક્ષીણતા (ડીજનરસી) & Degeneracy \\
ક્ષોભ સિદ્ધાંત & Perturbation theory \\
સમાન કણો & Identical particles \\
પૌલી બહિષ્કાર નિયમ & Pauli exclusion principle \\
ફર્મી ઊર્જા સ્તર & Fermi energy level \\
વસ્તી ઉલટાણ & Population inversion \\
પ્રેરિત ઉત્સર્જન & Stimulated emission \\
ક્વોન્ટમ અવરોધન & Quantum confinement \\
અવસ્થા-ઘનતા & Density of states \\
ક્યુબિટ & Qubit \\
અધિસ્થિતિ & Superposition \\
ગૂંથણ (જોડાણ) & Entanglement \\
ક્વોન્ટમ કી-વિતરણ & Quantum key distribution \\
\bottomrule
\end{longtable}

% ---------------------------------------------------------------------
\chapter[પરિશિષ્ટ E : સંદર્ભગ્રંથો]{સંદર્ભગ્રંથો અને વધુ વાંચન\\{\large (References and Further Reading)}}

\begin{enumerate}
  \item Beiser, A., \emph{Concepts of Modern Physics}, 6th ed., McGraw-Hill.
  \item Eisberg, R.\ \& Resnick, R., \emph{Quantum Physics of Atoms, Molecules, Solids,
        Nuclei and Particles}, 2nd ed., Wiley.
  \item Griffiths, D.\,J., \emph{Introduction to Quantum Mechanics}, 3rd ed., Cambridge
        University Press.
  \item Schiff, L.\,I., \emph{Quantum Mechanics}, 3rd ed., McGraw-Hill.
  \item Gasiorowicz, S., \emph{Quantum Physics}, 3rd ed., Wiley.
  \item Ghatak, A., \emph{Basic Quantum Mechanics}, Macmillan India.
  \item Mathur, D.\,S., \emph{Elements of Properties of Matter \& Atomic Physics}, Shukla
        \& Co. (ગુજરાત વિદ્યાર્થી-માર્ગદર્શન પરંપરા).
  \item NCERT/GSEB ધોરણ 11--12 ભૌતિકવિજ્ઞાન પાઠ્યપુસ્તકો, ગુજરાત રાજ્ય પાઠ્યપુસ્તક મંડળ,
        ગાંધીનગર.
  \item GTU Engineering Physics (3110016 / BE01000021) અભ્યાસક્રમ અને પ્રશ્ન-બેંક,
        ગુજરાત ટેકનોલોજિકલ યુનિવર્સિટી.
  \item ગુજરાત ગ્રંથ નિર્માણ બોર્ડ / યુનિવર્સિટી ગ્રંથ નિર્માણ બોર્ડ, ગુજરાત --
        ભૌતિકવિજ્ઞાન પરિભાષા-કોશ પરંપરા અને દ્વિભાષી પાઠ્ય-ગ્રંથ શ્રેણી.
  \item Nielsen, M.\,A.\ \& Chuang, I.\,L., \emph{Quantum Computation and Quantum
        Information}, Cambridge University Press.
  \item Planck, M. (1901), ``On the Law of Distribution of Energy in the Normal
        Spectrum'', \emph{Ann. der Physik}; Compton, A.\,H. (1923), \emph{Phys.\ Rev.},
        \textbf{21}, 483; Davisson--Germer (1927), \emph{Phys.\ Rev.}, \textbf{30}, 705.
\end{enumerate}

\vspace{6mm}
\begin{mulmudda}
\textbf{સમાપન}: આ ગ્રંથ ક્વોન્ટમ ભૌતિકવિજ્ઞાનના ઐતિહાસિક પ્રારંભથી ક્વોન્ટમ ગણતરી અને
ક્વોન્ટમ સંદેશાવ્યવહાર સુધીની સંપૂર્ણ સફર ગુજરાતી માધ્યમમાં આપે છે. વિદ્યાર્થીઓને સલાહ છે કે
દરેક ઉત્પત્તિ જાતે કરી જુઓ, TikZ આકૃતિઓના પરિમાણો બદલી પ્રયોગ કરો, અને MCQ/સ્વાધ્યાય
ઉકેલો સૂત્ર-સંગ્રહ (પરિશિષ્ટ B) સાથે સરખાવો.
\end{mulmudda}
